\documentclass[11pt,letterpaper]{article}

% Packages
\usepackage[margin=0.6in]{geometry}
\usepackage{enumitem}
\usepackage{hyperref}
\usepackage{titlesec}
\usepackage{xcolor}
\usepackage{amsmath,amssymb}

% Hyperref setup
\hypersetup{
    colorlinks=true,
    linkcolor=black,
    urlcolor=blue,
    citecolor=black
}

% Remove page numbers
\pagestyle{empty}

\setlength{\parindent}{0pt}

% Section formatting
\titleformat{\section}
    {\large\bfseries\uppercase}
    {}
    {0em}
    {}
    [\titlerule]

\titlespacing{\section}{0pt}{10pt}{5pt}

% Subsection formatting (for entries)
\titleformat{\subsection}
    {\bfseries}
    {}
    {0em}
    {}
    
\titlespacing{\subsection}{0pt}{6pt}{0pt}

% Custom commands
\newcommand{\institution}[1]{\textbf{#1}}
\newcommand{\position}[1]{\textit{#1}}
\newcommand{\dates}[1]{\hfill #1}

% Reduce spacing in lists
\setlist[itemize]{leftmargin=*,nosep,after=\vspace{3pt}}

\begin{document}

% Header
\begin{center}
    {\LARGE\textbf{Jefrey Bergl}} \\[3pt]
    \href{mailto:jbergl@unc.edu}{jbergl@unc.edu} $\mid$ 
    (336) 339-9069 $\mid$ 
    \href{https://qkdistributor.github.io/PersonalWebsite/}{Website} $\mid$
    \href{https://github.com/qkdistributor}{Github}
\end{center}

\vspace{-8pt}

% Education
\section{Education}

\institution{University of North Carolina at Chapel Hill} \dates{Expected Graduation: May 2026} \\
\position{\textbf{BS} Computer Science and Applied Mathematics} \\
Dean's List (5 semesters) \\
\textbf{Relevant Coursework:} Graduate Machine Learning, Graduate Computer Vision, Graduate Deep Learning, Optimization for Machine Learning, Numerical Analysis, Advanced Programming Language Concepts

% Research Experience
\section{Research Experience}

\institution{UNC Chapel Hill} \dates{Oct 2024 -- Present} \\
\position{Undergraduate Research Assistant} $\mid$ Advisor: Prof. Thomas Hofweber \\[-6pt]
\begin{itemize}
    \item Co-authored \textit{The Black Tuesday Attack: How to Crash the Stock Market With Adversarial Attacks to Financial Forecasting Models} (submitted to Journal of Cybersecurity; see Publications).
    \item Developed adversarial perturbation framework for financial forecasting models; ran experiments on CNN and MLP architectures.
    \item Received \$3,000 in undergraduate research funding to support the project.
\end{itemize}

% Publications and Preprints
\section{Publications \& Preprints}

Thomas Hofweber, \textbf{Jefrey Bergl}, Ian Reyes, Amir Sadovnik. \textit{The Black Tuesday Attack: How to Crash the Stock Market With Adversarial Attacks to Financial Forecasting Models}. Submitted to \textit{The Journal of Cybersecurity}, 2025. \href{https://arxiv.org/abs/2510.18990}{arXiv:2510.18990}.

% Projects
\section{Projects}

\textbf{H-JEPA: Learning Hamiltonian Dynamics on Latent Video Embeddings} [\href{https://github.com/qkdistributor/H-JEPA}{GitHub}] \hfill Dec 2025 \\
\textit{Independent Research} $\vert$ Python, PyTorch, V-JEPA 2 \hfill
\begin{itemize}
    \item Designed a latent world model that enforces Hamiltonian mechanics on frozen V-JEPA 2 video embeddings, learning a scalar energy function $H(q,p)$ and evolving states via a symplectic St\"{o}rmer--Verlet integrator with a learnable timestep
    \item Achieved 1.9$\times$ lower prediction MSE than a parameter-matched baseline MLP on 15-step rollouts, using 30\% fewer parameters (919K vs.\ 1.3M), with bounded energy drift by construction
    \item Conducted ablation studies isolating contributions of symplectic integration and energy conservation loss, revealing a prediction--conservation trade-off across four model variants
\end{itemize}

\textbf{Cost-Effective Brain Tumor Detection with CNNs} $\mid$ \textit{Python, TensorFlow/Keras, scikit-learn}
 \dates{Dec 2024} \\[-6pt]
\begin{itemize}
    \item Built low-parameter CNN to classify brain MRI scans, minimizing false negatives for prescreening in resource-limited settings.
    \item Compared against transfer learning baselines (ResNet-50, VGG-16); achieved Accuracy $\approx$ 0.85, F1 $\approx$ 0.87 while reducing false-negative rate and remaining far smaller to deploy. [\href{https://github.com/TrippWhaley/comp-562-final-project/}{Code}]
\end{itemize}

% Presentations
\section{Presentations}

\textbf{CoreAI UNC Research Showcase} -- ``The Black Tuesday Attack'' [\href{https://drive.google.com/file/d/1dxe09yK2t-jKYxDvMgamqlDRq_Eh0SYO/view?usp=sharing}{poster}]
 \dates{Feb 2026}

\textbf{AI@UNC} -- ``Adversarial Attacks on Financial Forecasting Models'' [\href{https://drive.google.com/file/d/1wtKcK2hbIyacfg8mf0JP0TJNtrU2I4Q9/view?usp=sharing}{slides}]
 \dates{April 2025}\\
\textit{Selected for Spotlight Presentation}

\textbf{COMP 790: Graduate Computer Vision} -- ``Evaluating Real World ViT Performance''
[\href{https://drive.google.com/file/d/1pDFAlfCgygF0m94B9xhhYrKzUEZAEjkT/view?usp=sharing}{poster}]
 \dates{May 2025}

% Technical Skills
\section{Technical Skills}

\textbf{Programming Languages:} Python, Java

\end{document}